\section{Support at the actual indices six through fourteen}
\label{ext:support6-14}
The support calculation at index five extends to each subsequent small
index, but its Gaussian expression must be evaluated at that index. We
record these evaluations explicitly; no numerical maximum at index five
is used at another index.

\begin{proposition}[Actual-index low-moment support]
\label{ext:support6-14-table}
Let $6\le n\le14$, assume $C_{n-1}\le\CE<C_n$, and let $(P,H)$ be
a positive closed-event ratio maximizer from
Corollary~\ref{ext:ratio-attainment}. For each applicable row below,
$\rho(P)$ in the stated closed interval implies
$\supp P\subset(-R,R)$. The number $b$ bounds
$|\beta|=|\E[X|X|]|$; the final column is a strict rational lower
bound for the endpoint margin, uniform over the indicated indices.
\begin{center}\small
\begin{tabular}{@{}ccccc@{}}
\toprule
indices & $\rho$ interval & $b$ & $R$ & margin $>$\\
\midrule
$6,\ldots,14$ & $[1,11/10]$ & $0.253257$ & $23/10$ & $0.259317$\\
$6,\ldots,14$ & $[11/10,23/20]$ & $0.307221$ & $12/5$ & $0.161008$\\
$6,\ldots,9$ & $[23/20,237/200]$ & $0.338920$ & $5/2$ & $0.182254$\\
$6,\ldots,9$ & $[23/20,6/5]$ & $0.351391$ & $13/5$ & $0.357463$\\
\bottomrule
\end{tabular}
\end{center}
The overlapping last two rows serve different purposes: the smaller
radius is used in the low-moment contact calculation, and the wider
moment interval is used by the supported Fourier estimate.
\end{proposition}

\begin{proof}
Put $C=C_n$, $s_n=\sqrt{n/(n-1)}$,
$e_n=n(s_n-1)-1/2$, and
$A=\sqrt{n-1}\Delta_{n-1}(P)/\rho$.
Then $d_P=n(A-C)\le0$. Define
\[
 G_n(H,t)=n\sqrt n\left\{
 \Phi\left(\frac{H-t}{\sqrt{n-1}}\right)
 -\Phi\left(\frac H{\sqrt n}\right)\right\}
 +nt\phi\left(\frac H{\sqrt n}\right)
 +\frac H2(t^2-1)\phi\left(\frac H{\sqrt n}\right).
\]
The contact inequality at the same marginal gives, at every support
point $t$,
\begin{equation}\label{ext:support6-14-contact}
 0\le G_n(H,t)
 +C\{-|t|^3+\tfrac32\rho t^2+3\beta t+\rho e_n\}
 +\rho s_n d_P.
\end{equation}
Indeed $n(s_nA-C)=s_nd_P+n(s_n-1)C$; together with
the constant $-C\rho/2$ in the contact polynomial this gives the
displayed $C\rho e_n$. Thus the immediate predecessor enters only
through the nonpositive last term.

For a table row let $r$ denote its upper moment endpoint and set
\[
 \begin{gathered}
 c=819/2000,\qquad E=91/1000,\\
 Q(v)=-v^3+\tfrac32rv^2+3bv+rE,\\
 L(v)=\tfrac25(v^2/2+1/2+E)+cQ'(v).
 \end{gathered}
\]
We have $c<\CE<C$ and $e_n<E$. The latter inequality can be
checked without a square root by
$((E+1/2+n)/n)^2>n/(n-1)$, for each of the nine indices.
The signed-moment inequality in Lemma~\ref{ms:sign-moment} and the
strict rational comparison $(1+b^2)^2>r^2(1-b^2)$ give
$|\beta|\le b$.
Every row satisfies
\[
 Q(R),Q'(R),Q''(R),L(R),L'(R)<0,\qquad L''=2/5-6c<0.
\]
Hence $Q(v)<0$ for $v\ge R$, and the polynomial term in
\eqref{ext:support6-14-contact} is at most $cQ(|t|)$.

The Gaussian derivative estimate used at index five holds with the
actual index:
\begin{equation}\label{ext:support6-14-derivative}
 |\partial_tG_n(H,t)|
 \le\phi(0)\{t^2/2+n(s_n-1)\}
 <\tfrac25(t^2/2+1/2+E).
\end{equation}
For completeness, with $z=H/\sqrt n$ and $y=(H-t)/\sqrt n$,
the derivative is $n\phi(z)+Ht\phi(z)-ns_n\phi(s_ny)$.
Taylor's theorem for $\phi$ bounds
$|n\phi(y)-n\phi(z)-Ht\phi(z)|$ by $\phi(0)t^2/2$.
The identity $\partial_s[s\phi(sy)]=-\phi''(sy)$ bounds
the remaining scaling error by $n\phi(0)(s_n-1)$.
For either sign $\eta$, the derivative of
$G_n(H,\eta v)+cQ(v)$ is therefore at most $L(v)$, which is
negative on $[R,\infty)$. It remains to exclude $v=R$ for every
real threshold $H$.

Differentiating twice in $H$ gives the uniform bound
\[
 |\partial_H^2G_n(H,R)|\le
 M_{n,R}:=\frac{\sqrt n}{4}\left(\frac n{n-1}+1\right)
 +\frac25R+\frac{2|R^2-1|}{5\sqrt n}.
\]
Here one uses $\|\phi'\|_\infty<1/4$,
$\|\phi''\|_\infty<2/5$, and
$\|\phi'''\|_\infty<4/5$. Set $H_0=35/2$ for
$6\le n\le9$, and $H_0=22$ for $10\le n\le14$.
The rational interval evaluations at the nodes
$-H_0,-H_0+1/10,\ldots,H_0$, enlarged by $M_{n,R}/800$,
cover every closed interval between consecutive nodes. This is the
standard second-derivative interpolation error $Mh^2/8$, applied to
both signs of $G_n$.

For both exterior rays, the exact squared comparisons
$H_0^2\ge25n$ and $(H_0-R)^2\ge25(n-1)$ put the two
Gaussian arguments in the same tail beyond five standard deviations.
Consequently, with $\tau=10^{-6}>\Phi(-5)$,
\[
 |G_n(H,R)|\le n\sqrt n\,\tau+
 \left(nR+\frac{H_0}2|R^2-1|\right)
 \phi(0)e^{-H_0^2/(2n)},\qquad |H|\ge H_0.
\]
The difference of the two distribution functions is bounded by
$\tau$, since both lie in the same tail interval of length $\tau$.
Also $|H|\phi(H/\sqrt n)$ decreases on these rays.
The reflection identity $G_n(H,-R)=-G_n(-H,R)$ supplies the other
support endpoint.

Let $B_{n,R}$ be the larger of the central interpolation bound and
the exterior bound. Exact rational arithmetic proves
$B_{n,R}+cQ(R)<0$ in all 26 rows; the table gives lower bounds for
$-cQ(R)-B_{n,R}$. Ray monotonicity now contradicts
\eqref{ext:support6-14-contact} whenever $|t|\ge R$.
\end{proof}

\paragraph{Reproduction of the table.}
The independent programs \texttt{check\_support6\_14.py} and
\texttt{check\_support6\_9\_bc.py} use the rational Gaussian,
exponential and square-root enclosures described after
Proposition~\ref{ext:n5-support-table}. They check respectively
8,600 and 1,400 closed threshold cells, both exterior rays, and all
moment and derivative gates. The actual integer $n$ occurs in every
Gaussian expression. The saved records retain the full rational
margins; each displayed decimal is an exact rational rounded down.
These support conclusions apply to violating ratio maximizers under
the stated predecessor hypothesis. They do not assert that arbitrary
laws with a bounded third moment have bounded support.


\begin{proposition}[Support for the complementary Fourier intervals]
\label{ext:support6-14-high-table}
Under the same predecessor and ratio-maximizer hypotheses, each row below
also gives $\supp P\subset(-R,R)$ on its whole stated moment interval.
The signed-moment bound and endpoint-margin convention are those of
Proposition~\ref{ext:support6-14-table}.
\begin{center}\small
\begin{tabular}{@{}ccccc@{}}
\toprule
indices & $\rho$ interval & $b$ & $R$ & margin $>$\\
\midrule
$6,\ldots,9$ & $[6/5,7/5]$ & $0.478710$ & $3$ & $0.281666$\\
$6,\ldots,9$ & $[7/5,3/2]$ & $0.525521$ & $16/5$ & $0.316642$\\
$6,\ldots,9$ & $[3/2,8/5]$ & $0.565414$ & $17/5$ & $0.392665$\\
$6,\ldots,9$ & $[8/5,17/10]$ & $0.600001$ & $18/5$ & $0.516295$\\
$6,\ldots,9$ & $[17/10,7361/4000]$ & $0.641563$ & $19/5$ & $0.278368$\\
$10,\ldots,14$ & $[23/20,71/50]$ & $0.488731$ & $3$ & $0.139019$\\
\bottomrule
\end{tabular}
\end{center}
\end{proposition}
\begin{proof}
Use the same $c=819/2000$, $E=91/1000$, contact inequality and
polynomials $Q,L$ as in the preceding proof, substituting each listed
$(r,b,R)$. Every signed-moment and derivative gate holds strictly.
For these larger support radii take
\[
 H_0=35/2\quad(6\le n\le8),\qquad
 H_0=18\quad(n=9),\qquad H_0=22\quad(10\le n\le14).
\]
These choices satisfy both squared tail comparisons in that proof.
The enlargement at $n=9$ is needed for $R=19/5$:
$(18-19/5)^2\ge25(9-1)$, whereas $H_0=35/2$ would not satisfy
this condition. Apply the same interpolation bound on the closed
mesh of width $1/10$ and the same estimate on both exterior rays.

The program \texttt{check\_support6\_14\_high.py} evaluates all
$25$ actual-index rows, comprising $9250$ closed threshold cells.
Its exact rational margins are strictly larger than the displayed
rounded-down values. Thus $B_{n,R}+cQ(R)<0$ throughout each row;
ray monotonicity excludes every support point with $|t|\ge R$.
These bounds supply the supported Fourier premises in the
complementary-moment part of Section~\ref{ext:finite-five-fourteen}.
\end{proof}
